\section{Multi-Seed Robustness}

\subsection{Scenario outcomes}

\begin{table}[H]
\centering
\footnotesize
\input{build/technical/v5.0.8/includes/table_technical_robustness}
\caption{Technical outcome, fairness disposition, headline gender Female/Male ratio, and the
lowest decision-driving comparison over ten fixed seeds. Each driver is read from the run's
\texttt{compliance\_decision.basis.worst\_case}; the tested group or intersection, reference,
ratio, threshold status, and seed stay together. A technical pass and a fairness review warning are
separate fields.}
\label{tab:robustness}
\end{table}

All \RobustnessPasses{} configured scenario runs passed their technical contracts. Balanced and
security scenarios keep the headline disparity ratio close to one. The outlier scenario spans
\OutliersRatioMin{} to \OutliersRatioMax{}. The gender-bias positive control consistently produces a
ratio near 0.65: range \GenderBiasRatioMin{}--\GenderBiasRatioMax{}, mean
\GenderBiasRatioMean{}. Every one of those ten runs is outside the fairness screen, as expected by
the test design, while still passing the scenario contract.

This distinction matters for evaluators. The positive control demonstrates that the pipeline does
not erase a planted disparity and that the review status travels with the controlling comparison.
Across its plan, the lowest gender-bias decision driver is the Female--Asian intersection versus
the evidence-recorded male--White intersection at 0.532605 (seed 101). For seed 42, the driver is
Female--White versus male--White with ratio 0.611312. In the balanced plan, the lowest driver is the
race Other group versus the observed-rate-selected Hispanic reference at 0.887639 (seed 101).
Section~\ref{sec:pair-selection} explains this release's observed-rate race-reference policy and its
post-selection limitation. The public 40-row projection keeps scenario, seed, controlling
attribute, tested group or intersection, reference, ratio, threshold and status, technical scenario
status, AOD/EOD availability, and integrity results together. It retains the lowest intersectional
comparison as a separately named diagnostic rather than treating it as the decision driver.

\subsection{Evidence integrity across the plan}

Every projected row records successful PDF numeric parity, certificate-chain validation,
certificate-signature validation, evidence encryption, and internal LC11 contract-certificate
validation. The public
projection is sanitized: it excludes private runner paths, logs, per-seed working directories,
row-level data, and the full private Gold tree.

A 40/40 result is evidence that this exact planned set completed consistently. It is not a confidence
interval over all future datasets, configurations, operating systems, dependency versions, or
adversarial conditions. The seed list and four scenarios define the supported robustness claim.

\subsection{Advisory utility limitation}

Amplification utility AUC ranges from \UtilityMinAUC{} to \UtilityMaxAUC{}, with mean
\UtilityMeanAUC{}. Only \UtilityAtOrAboveFloor{} of 40 runs are at or above the configured 0.70
floor; \UtilityBelowFloor{} are warnings. Accuracy ranges from \UtilityMinAccuracy{} to
\UtilityMaxAccuracy{}, with mean \UtilityMeanAccuracy{}; no accuracy threshold is configured, so
these values have no pass/warn disposition. The result is stable enough to expose a consistent
limitation: the present public evidence does not establish strong train-on-synthetic predictive
utility even on these generated input fixtures. No intrinsic utility result exists.

The utility diagnostic and the quality composite answer different questions. Distribution,
correlation, demographic-rate fidelity, and a privacy heuristic can yield a composite above 0.80
while the advisory AUC remains below 0.70. The paper therefore reports both and does not use the
quality gate to imply predictive performance.
